\abstractzh{
新兴经济体资本市场自由化水平不足不利于企业融资，进而制约其转型发展。本文以中国A股分阶段纳入MSCI新兴市场指数为外生制度冲击，基于2010-2022年间2,712家上市公司数据，借助双重差分识别策略，并结合事件研究和倾向得分匹配等方法，研究新兴市场经济体的资本市场自由化能否推动企业数字化转型。

实证结果显示，在TWFE基准设定下，MSCI纳入对应的数字化转型系数为0.39个对数点（约48\%）。但在分期处理与队列ATT等现代DID诊断中，平均效应明显收敛且统计精度下降，表明传统TWFE在异质处理效应下可能高估总体效应。事件研究结果显示，处理前不存在系统性趋势差异，处理后效应呈累积特征。安慰剂检验、替代标准误聚类、受限样本与替代被解释变量检验在方向上支持主结论。异质性分析发现，效果主要集中于高技术行业和制度环境较好的地区，说明吸收能力与制度互补性是关键约束条件。

机制检验显示，融资渠道证据具有“结构调整强、综合约束弱”的特征：银行杠杆和现金流指标显著改善，但KZ指数不显著；治理渠道指标方向一致且显著，但需考虑合规性与被动持股稀释等制度解释。知识溢出渠道因关键中介变量未完全并入样本，暂不作直接系数结论。总体而言，国际金融一体化对企业数字化升级具有条件性促进作用，但效应大小依赖估计方法、测度设定与企业实施能力。
政策含义上，资本市场开放改革需要与信息披露质量提升、公司治理改进和区域数字基础条件协同推进，才能更稳定地转化为企业持续数字化投入。对企业而言，纳入带来的融资机会与外部监督约束并存，长期收益仍取决于组织执行能力与技术吸收能力。
}
{
资本市场自由化；数字化转型；MSCI纳入；中国A股；新兴市场；计算文本分析；双重差分
}

\abstracten{
The insufficient level of capital market liberalization in emerging economies hinders corporate financing, thus restricting their transformation and development. This paper takes China's phased inclusion in the MSCI Emerging Markets Index as an exogenous institutional shock. Based on data from 2,712 listed companies between 2010 and 2022, and using a difference-in-differences (DID) identification strategy, along with event studies and propensity score matching, this study investigates whether capital market liberalization in emerging market economies can promote corporate digital transformation.

Empirical results show that in the TWFE benchmark specification, MSCI inclusion is associated with a 0.39 log-point increase in digital transformation intensity (about 48\%). However, modern staggered DID diagnostics based on cohort-specific ATT estimates are materially smaller and less precise, indicating substantial estimator sensitivity under heterogeneous treatment timing. Event-study evidence supports no systematic pre-treatment divergence and shows cumulative post-treatment dynamics. Placebo tests, alternative clustering strategies, restricted-sample checks, and alternative outcomes provide directional robustness.

Heterogeneity results indicate stronger effects in high-technology sectors and institutionally developed regions, consistent with an absorptive-capacity interpretation. Mechanism evidence is nuanced: financing-channel results are more consistent with financing-structure adjustment than broad constraint relaxation (bank leverage and cash-flow proxies are significant, while the KZ index is not), governance proxies are directionally strong but institutionally interpretable with compliance and passive-ownership caveats, and direct knowledge-spillover estimation remains deferred due to mediator-data limitations. Overall, international financial integration appears to be a conditional contributor to digital upgrading in emerging markets, with effect magnitudes dependent on estimator choice, measurement assumptions, and firm-level implementation capacity.
From a policy perspective, market-opening reforms should be coordinated with stronger disclosure quality, governance discipline, and regional digital support conditions to translate external capital exposure into sustained firm-level digital investment. For firms, inclusion-related financing opportunities and external monitoring pressures coexist, and long-run gains depend on implementation and absorptive capacity rather than inclusion status alone.
}
{
Capital market liberalization; Digital transformation; MSCI inclusion; Chinese A-shares; Emerging markets; Computational text analysis; difference-in-differences
}
